\documentclass[11pt]{article}
\usepackage[margin=1in]{geometry}
\usepackage{amsmath,amssymb,amsthm}
\usepackage{booktabs}
\usepackage{enumitem}
\usepackage{float}
\usepackage{array}
\usepackage{xcolor}
\usepackage{tikz}
\usepackage{xurl}
\usepackage[hidelinks]{hyperref}

\newtheorem{theorem}{Theorem}
\newtheorem{lemma}[theorem]{Lemma}
\newtheorem{proposition}[theorem]{Proposition}
\newtheorem{corollary}[theorem]{Corollary}
\newcommand{\cube}[1]{Q_{#1}}
\newcommand{\betaQ}[1]{\beta(\cube{#1})}
\newcommand{\bits}{\{0,1\}}
\newcommand{\signs}{\{-1,+1\}}
\newcommand{\ternary}{\{-1,0,1\}}
% Author-facing note flagged for removal before final. Renders loudly on purpose.
\newcommand{\TODO}[1]{\textcolor{red}{\textbf{[TODO --- cut before final:]}~#1}}

\title{Exact Metric Dimensions and Improved Upper Bounds for Hypercubes}
\author{Shmulik Cohen}
\date{August 2026}

\begin{document}
\maketitle

\begin{abstract}
In this paper we consider the problem of determining the metric dimension of
the hypercube \(Q_n\). Exact values are difficult to obtain even for modest
dimensions, and the previously established sequence stopped at \(n=13\). We
prove that \(\beta(Q_{14})=\beta(Q_{15})=9\). The lower bound is obtained by an
exhaustive enumeration of detecting matrices, using hypercube symmetries and
orderly generation to reduce the search. The computation is implemented
separately in C and safe Rust and is supported by direct checks on smaller
instances and comparisons with independent data. As a consequence we also
determine the nonadaptive coin-weighing value \(M(14)=8\). Finally, we give
explicit, independently checkable matrices that improve the known upper bounds
to \(\beta(Q_{24})\leq13\), \(\beta(Q_{26})\leq14\), and
\(\beta(Q_{29})\leq15\).
\end{abstract}

\section{Introduction}

The vertices of the hypercube \(Q_n\) are the binary vectors \(\bits^n\), with
two vertices adjacent when they differ in one coordinate. A set \(R\) of
vertices is \emph{resolving} if the vectors of Hamming distances
\[
  \bigl(d(x,r):r\in R\bigr), \qquad x\in\bits^n,
\]
are pairwise distinct. The minimum cardinality of a resolving set is the
metric dimension \(\beta(Q_n)\).

The metric-dimension problem was introduced independently by Slater
\cite{slater1975} and Harary and Melter \cite{harary1976}. For hypercubes it is
closely related to nonadaptive coin weighing and detecting matrices
\cite{erdosrenyi1963,cantormills1966,sebotannier2004,luye2022}. Its asymptotic
order is known, but exact small values are difficult. The public sequence
A303735 records exact values through \(n=13\) as
\[
1,2,3,4,4,5,6,6,7,7,8,8,8
\]
\cite{oeisA303735}. The smaller terms have since been independently confirmed:
V.~Miller verified the first ten values with a MaxSat solver and
symmetry-breaking constraints \cite{oeisA303735}. Heuristic searches had already
produced the upper bounds
\(\beta(Q_{14})\leq9\) and \(\beta(Q_{15})\leq9\)
\cite{mladenovic2012,nikolic2017,hertz2020}. What remained open was a lower
bound ruling out eight landmarks.

This paper supplies that lower bound by exhaustive computation. The main
result is the following.

\begin{theorem}\label{thm:main}
The following statements hold.
\begin{enumerate}[label=(\roman*)]
\item The metric dimensions of the 14- and 15-dimensional hypercubes are
\[
  \beta(Q_{14})=\beta(Q_{15})=9.
\]
\item The metric dimensions of the 24-, 26-, and 29-dimensional hypercubes
satisfy
\[
  \beta(Q_{24})\leq13,\qquad
  \beta(Q_{26})\leq14,\qquad
  \beta(Q_{29})\leq15.
\]
\end{enumerate}
\end{theorem}

The computation uses familiar hypercube symmetries. Nikoli\'c et al. showed,
among other reductions, that a resolving set may be
translated to contain the zero vertex and that complementing landmarks
preserves resolving pairs \cite{nikolic2017}. Their implementation used these
facts in a greedy upper-bound heuristic and used lexicographic order for
candidate tie-breaking. Our lower-bound method differs in kind: we combine the
normalized detecting-matrix formulation with an exhaustive sorted-prefix
search, rejecting a prefix whenever it is not the lexicographically least
member of its symmetry orbit. The prefix rule is proved in
Section~\ref{sec:orderly}; it is established directly, not inferred from
heuristic success.

The code, certificates, complete correctness argument, and reproduction scripts
are available in the accompanying repository \cite{artifact}. The exact claim
does not depend on the stochastic searches used to discover the upper bounds.

\section{Resolving sets and detecting matrices}\label{sec:reduction}

Let \(W\in\signs^{q\times n}\). We call \(W\) a \emph{detecting matrix} if
\[
 Wz\neq0 \quad\text{for every nonzero }z\in\ternary^n.
\]

\begin{lemma}\label{lem:equivalence}
There is a resolving set of size \(q\) in \(Q_n\) if and only if there is a
\(q\times n\) detecting matrix.
\end{lemma}

\begin{proof}
Given landmarks \(v_1,\ldots,v_q\in\bits^n\), set
\(W_{ij}=1-2(v_i)_j\). For \(x,y\in\bits^n\), direct expansion gives
\[
 d(x,v_i)-d(y,v_i)=\sum_{j=1}^n (x_j-y_j)W_{ij}.
\]
Thus two distinct vertices have identical distance vectors precisely when the
nonzero vector \(z=x-y\in\ternary^n\) lies in the kernel of \(W\). Conversely,
every nonzero \(z\in\ternary^n\) can be written as \(x-y\) for binary
\(x\neq y\), so the implication reverses.
\end{proof}

Multiplying a column of \(W\) by \(-1\) preserves the condition: if
\(D\) is a diagonal sign matrix, then \(WDz=0\) if and only if
\(W(Dz)=0\), and \(D\) permutes \(\ternary^n\). We may therefore flip
columns so that the first row is all \(+1\).

\begin{lemma}[subset-sum form]\label{lem:subsetsums}
A sign matrix is detecting if and only if the \(2^n\) subset sums of its
columns are pairwise distinct.
\end{lemma}

\begin{proof}
Two equal subset sums, indexed by \(A,B\subseteq[n]\), give a ternary kernel
vector after canceling \(A\cap B\): assign \(+1\) to \(A\setminus B\),
\(-1\) to \(B\setminus A\), and zero elsewhere. Conversely, the positive
and negative supports of any ternary kernel vector give two equal subset sums.
\end{proof}

In the language of additive combinatorics, Lemma~\ref{lem:subsetsums} says that
the columns of a detecting matrix form a \emph{dissociated} set: a set of
vectors with no nontrivial \(\ternary\)-combination summing to zero,
equivalently a set whose subset sums are pairwise distinct
\cite{mathoverflow157634}. The metric dimension \(\beta(Q_n)\) is thus the least
number of coordinates \(q\) in which \(n\) sign vectors can be dissociated;
Theorem~\ref{thm:main}(i) states that no fourteen vectors in \(\signs^8\) are
dissociated, whereas fifteen can be dissociated in nine coordinates. Because the
first row is normalized to all \(+1\), any two equal subset sums must have equal
cardinality; under the coordinate bijection \(u\mapsto\frac{1-u}{2}\), collisions
in \(\signs^{q-1}\) between subsets of the same size map bijectively to
\(\bits^{q-1}\) subset-sum collisions. In this \(\bits\)-valued form, maximal
systems are catalogued as maximum dissociated sets of \(0\)--\(1\) vectors
\cite{mckaydata}, which is why those catalogues furnish an independent check in
Section~\ref{sec:verification}.

After normalization, a column is determined by its last \(q-1\) signs. We
encode it as a point of \(\bits^{q-1}\), so there are only \(2^{q-1}\)
column types. Repeated columns are impossible, since their difference is a
nonzero ternary kernel vector. The search object is therefore an \(n\)-subset
of a \(2^{q-1}\)-element type set.

A small example illustrates the role of the normalized row. For four columns,
write
\[
W=\left[
\begin{array}{c|rrrr}
 &c_1&c_2&c_3&c_4\\ \hline
\text{normalized row}&+1&+1&+1&+1\\
 &+1&-1&+1&-1\\
 &-1&+1&+1&-1
\end{array}\right].
\]
The first coordinate of a subset sum is its cardinality. Hence subsets of
different sizes cannot collide, while the remaining rows must distinguish
subsets of the same size. Normalization does not itself establish the detecting property;
it exposes a canonical coordinate and reduces each column to one of
\(2^{q-1}\) types.

\section{Incremental exhaustive search}\label{sec:dfs}

Represent a partial object by an increasing list of types
\(P=(c_1<\cdots<c_k)\). Let
\[
 \Sigma(P)=\left\{\sum_{c\in A}c:A\subseteq P\right\}
\]
denote its subset sums in the original sign-vector coordinates. If all sums
in \(\Sigma(P)\) are distinct, adjoining a candidate \(c\) is legal exactly
when
\[
 \Sigma(P)\cap\bigl(\Sigma(P)+c\bigr)=\varnothing.       \tag{1}
\]
The child sums are then the disjoint union of these two sets. The enumerator
stores the sums once per depth and tests (1) with a hash table.

Without symmetry pruning, the depth-first traversal is complete: at depth
\(k\), every legal larger candidate is tried in increasing order and only the
suffix after that candidate is passed to the child. Thus every collision-free
sorted type set containing the fixed first type appears once.

Two elementary pruning rules are monotone.
\begin{enumerate}[label=(\roman*)]
\item If a candidate \(c\) is illegal for \(P\), equation (1) remains false
for every extension of \(P\), because the witnessing subset sums remain
present.
\item If fewer than \(n-k\) live candidates remain, the prefix cannot reach
\(n\) columns.
\end{enumerate}

For the decisive \(8\times14\) case, each coordinate of a subset sum is packed
into one byte of a 64-bit word with bias 64. With at most 14 columns every
coordinate lies between 50 and 78, so no carry or borrow can cross a byte
boundary. The packed addition is therefore exact componentwise arithmetic,
not a probabilistic hash.

\section{Orderly generation and prefix pruning}\label{sec:orderly}

The full automorphism group acting on $q\times n$ sign matrices is the signed
permutation group $O(q, \mathbb{Z}) = (\mathbb{Z}_2)^q \rtimes S_q$ on the rows
(of order $2^q q!$), combined with column permutations $S_n$ and column sign flips
$(\mathbb{Z}_2)^n$. Fixing the first row to all $+1$ by column sign flips reduces
each column to one of $2^{q-1}$ normalized column types $V = \{w \in \signs^q :
w_1 = +1\} \cong \bits^{q-1}$.

For any hypercube isometry $g \in O(q, \mathbb{Z})$ acting on $\signs^q$, we define
its normalized action $\Phi_g: V \to V$ by
\[
 \Phi_g(w) = (g(w))_1 \cdot g(w).
\]
Because $(g(w))_1 \in \signs$, the first coordinate of $\Phi_g(w)$ is identically
$+1$. The map $\Phi_g$ is a pointwise bijection on the finite set $V$: if
$\Phi_g(u) = \Phi_g(v)$, then $g(u) = \sigma g(v) = g(\sigma v)$ for
$\sigma = (g(u))_1(g(v))_1 \in \signs$; invertibility of $g$ gives $u = \sigma v$,
and comparing first coordinates yields $u_1 = +1 = \sigma v_1 = \sigma$, so $u = v$.

For group elements fixing the first row, $g \in \mathrm{Stab}(\text{row 1}) \cong
S_{q-1} \ltimes (\mathbb{Z}_2)^{q-1}$, the action on $\bits^{q-1}$ reduces to
\[
 x\longmapsto \pi(x\oplus t),
\]
where $\pi\in S_{q-1}$ permutes the lower coordinates and $t\in\bits^{q-1}$ is a
translation. For elements that pivot a different row $R \in \{1,\dots,q-1\}$ to
row 1, $\Phi_g$ acts on the type bits $(b_0,\dots,b_{q-2}) \in \mathbb{F}_2^{q-1}$
as the invertible linear map
\[
 \text{bit}_{\text{new Row 0}}(t) = b_{R-1}, \qquad
 \text{bit}_{\text{new Row } r'}(t) = b_{R-1} \oplus b_{r'-1}.
\]

We sort every type set lexicographically. Translation allows one chosen member of
a set to become zero, so every orbit has a sorted representative containing zero;
the search fixes zero as its first type. At a partial set $P$, the program tests
the normalized group actions and rejects $P$ if a sorted image is
lexicographically smaller.

\begin{proposition}[canonical-prefix rule]\label{prop:prefix}
If a sorted full set $X \subset V$ is lexicographically least in its orbit under
the normalized group action, then every initial prefix $A \subseteq X$ is
lexicographically least among its normalized sorted images $\Phi_g(A)$ that
contain zero. Consequently, rejecting a nonminimal prefix cannot remove the last
representative of a solution orbit.
\end{proposition}

\begin{proof}
Let $A$ be the first $k$ elements of $X$. Because $\Phi_g: V \to V$ is a pointwise
bijection, $\Phi_g(A) \subseteq \Phi_g(X)$. Let sorted $\Phi_g(A) = (y_1 < \dots < y_k)$
and sorted $\Phi_g(X) = (z_1 < \dots < z_n)$. Since $\Phi_g(A)$ is a $k$-subset of
$\Phi_g(X)$, the $j$-th smallest element of the superset is at most the $j$-th
smallest element of the subset: $z_j \le y_j$ for all $1 \le j \le k$.

Suppose some group element $g$ yields a normalized sorted image $\Phi_g(A)$ strictly
lexicographically smaller than $A$. Let $m \in \{1,\dots,k\}$ be the first index
where $y_m < x_m$ (with $y_j = x_j$ for $j < m$). For all $j < m$, $z_j \le y_j = x_j$.
If $z_j < x_j$ for any $j < m$, then sorted $\Phi_g(X) <_{\mathrm{lex}} X$. If
$z_j = x_j$ for all $j < m$, then at index $m$, $z_m \le y_m < x_m$, which again
yields sorted $\Phi_g(X) <_{\mathrm{lex}} X$. In either case, $X$ was not
lexicographically minimal in its orbit, a contradiction.
\end{proof}

This is an instance of Read--Farad\v zev orderly generation
\cite{read1978,faradzev1978}. For $q=8$, the stabilizer group has order
$7!\,2^7=645{,}120$, while the full row group has order $8!\,2^8 = 10{,}321{,}920$.
The baseline enumeration canonicalizes under the stabilizer through depth eight,
which is already sufficient to complete the search deterministically. Stopping early
may repeat work but cannot discard a solution.

The minimality test itself avoids running all \((q-1)!\) permutations on most
prefixes. A permutation only reorders coordinates, so it preserves the popcount
of every column type; the least type index attainable for a type of popcount
\(w\) is therefore fixed in advance. Comparing these popcount-derived bounds
against the incumbent already settles the ordering for the large majority of
prefixes, and the full permutation sweep is run only for the few that survive
this cheap filter.

\begin{figure}[t]
\centering
\begin{tikzpicture}[
  x=1.0cm, y=1.15cm,
  kept/.style={circle, fill=black, inner sep=1.5pt},
  pruned/.style={circle, draw=gray, dashed, inner sep=1.5pt},
  ke/.style={thick},
  pe/.style={dashed, gray},
  rlbl/.style={font=\footnotesize\itshape, gray},
  dep/.style={font=\footnotesize, anchor=east},
]
% left depth/count axis
\node[dep] at (-4.2,0)    {depth 1: \(1\)};
\node[dep] at (-4.2,-1.5) {depth 2: \(6\)};
\node[dep] at (-4.2,-3.0) {depth 3: \(18\)};
\node[dep] at (-4.2,-5.3) {depth 6: \(11{,}364\)};
% root
\node[kept] (r) at (0,0) {};
\node[font=\footnotesize, anchor=south] at (0,0.18) {\(c_1=0\) fixed};
% depth 2
\node[kept]   (a) at (-2.2,-1.5) {};
\node[kept]   (b) at (-0.7,-1.5) {};
\node[kept]   (c) at ( 0.8,-1.5) {};
\node[pruned] (d) at ( 2.3,-1.5) {};
\draw[ke] (r)--(a); \draw[ke] (r)--(b); \draw[ke] (r)--(c); \draw[pe] (r)--(d);
\node[rlbl, anchor=west] at (2.55,-1.5) {non-canonical};
% depth 3 (expand b)
\node[kept]   (e) at (-1.6,-3.0) {};
\node[pruned] (f) at (-0.7,-3.0) {};
\node[kept]   (g) at ( 0.2,-3.0) {};
\draw[ke] (b)--(e); \draw[pe] (b)--(f); \draw[ke] (b)--(g);
\node[rlbl, anchor=north] at (-0.7,-3.25) {collision, eq.~(1)};
\node[font=\footnotesize] at (-2.2,-2.35) {\(\vdots\)};
\node[font=\footnotesize] at ( 0.8,-2.35) {\(\vdots\)};
\node at (0,-3.95) {\(\vdots\)};
% frontier
\draw[dashed, thick] (-4.0,-5.3) -- (4.0,-5.3);
\node[font=\footnotesize, anchor=south] at (0,-5.22) {depth-6 frontier};
% partitions
\draw (-2.6,-5.3)--(-3.0,-6.2)--(-2.2,-6.2)--cycle;
\node[font=\scriptsize,anchor=north] at (-2.6,-6.2) {\(P_0\)};
\draw (-1.3,-5.3)--(-1.7,-6.2)--(-0.9,-6.2)--cycle;
\node[font=\scriptsize,anchor=north] at (-1.3,-6.2) {\(P_1\)};
\node[font=\footnotesize] at (0.0,-5.75) {\(\cdots\)};
\draw (1.3,-5.3)--(0.9,-6.2)--(1.7,-6.2)--cycle;
\node[font=\scriptsize,anchor=north] at (1.3,-6.2) {\(P_9\)};
\end{tikzpicture}
\caption{Schematic of the canonical-prefix enumeration for the \(8\times14\)
case. Each node is a sorted, collision-free prefix; depth is the number of
columns selected, with the first type fixed to zero. Solid edges are explored;
a dashed edge is pruned, either because the candidate would create a subset-sum
collision (equation~(1)) or because the prefix is not lexicographically least in
its symmetry orbit (Proposition~\ref{prop:prefix}). The counts on the left are
the numbers of canonical prefixes at each depth. The depth-six frontier of
\(11{,}364\) prefixes is split into ten congruence classes \(P_0,\ldots,P_9\)
modulo~\(10\), each searched independently and each reporting no \(8\times14\)
detecting matrix.}
\label{fig:tree}
\end{figure}

Figure~\ref{fig:tree} sketches the pruned enumeration and the depth-six frontier
at which the computation is partitioned; the scale of the reduction is shown in
Table~\ref{tab:pruning}. These counts
come from a controlled \(7\times11\) run and from the common
\(8\times14\) frontier; they describe performance, not an additional premise
of correctness.

\begin{table}[H]
\centering
\caption{Effect of canonical-prefix pruning.}
\label{tab:pruning}
\begin{tabular}{lrr}
\toprule
Measurement & Without prefix test & With prefix test\\
\midrule
\(7\times11\), complete search nodes & 256,123,828 & 189,831\\
\(8\times14\), nodes at depth 5 & 1,282,411 & 1,082\\
\(8\times14\), nodes at depth 6 & 34,942,921 & 11,364\\
\bottomrule
\end{tabular}
\end{table}

Here \emph{depth} means the number of columns already selected. Thus depth
six is the frontier of all legal canonical six-column prefixes, not six levels
of a binary decision tree.

\section{Partitioning the decisive computation}\label{sec:partition}

The \(8\times14\) traversal is divided into ten processes at depth six. Every
process deterministically walks the same tree down to that frontier. For each
outgoing child it increments the same counter; partition \(p\) retains the
child exactly when the counter is congruent to \(p\pmod {10}\). Every outgoing
branch therefore belongs to exactly one partition.

All completed logs have the identical frontier vector
\[
 (N_0,\ldots,N_6)=(0,1,6,18,143,1082,11364).
\]
The ten reported totals sum to 340,092,619 node visits. The common 12,614
shallow nodes (\(\sum_{i=0}^6 N_i = 12{,}614\)) are visited by every process,
so nine duplicate copies must be removed. The number of distinct search-tree
nodes is therefore
\[
 340{,}092{,}619-9\cdot12{,}614=339{,}979{,}093.
\]
Each partition exits successfully and reports zero solutions.

The partition-log checker is conservative by design: it accepts a run only when
the logs comprise exactly partitions 0 through 9, exactly one completed verdict
per file, the common frontier above, and no emitted solution. Malformed,
duplicate, missing, or interrupted logs are rejected.

\section{Verification and computational results}\label{sec:verification}

The proof obligation is narrower than the observation that two independent
programs reported no solution. It is that the implemented traversal covers every
normalized object up to the proved group action, and that every accepted pruning
rule preserves at least one representative. Sections~\ref{sec:reduction}--\ref{sec:partition} establish
those facts. We used the following controls to test the implementation against
those arguments.

\begin{table}[H]
\centering
\caption{Verification map.}
\label{tab:verification}
\begin{tabular}{>{\raggedright\arraybackslash}p{0.25\linewidth}>{\raggedright\arraybackslash}p{0.68\linewidth}}
\toprule
Control & Result\\
\midrule
Direct Python enumeration & Exact raw solution sets agree on \(3\times4\),
\(4\times4\), and \(5\times6\).\\
Whole-orbit canonicalization & All 13 independently computed canonical
\(5\times6\) representatives agree.\\
Partition control & Three small-instance parts are disjoint and their union is
the direct reference set.\\
Runtime instrumentation & AddressSanitizer and UndefinedBehaviorSanitizer pass
the small complete comparisons.\\
Known-value ladder & The code reproduces
\(\beta(Q_{11})=\beta(Q_{12})=\beta(Q_{13})=8\).\\
Independent data & All 118,485 McKay dissociated 12-set orbits occur in the
\(8\times13\) output after the proved conversion; none is missing
\cite{mckaydata}.\\
Second implementation & Safe Rust, with no external crates and no
\texttt{unsafe}, matches every complete C node vector in all ten
\(8\times14\) partitions.\\
\bottomrule
\end{tabular}
\end{table}

On the same arm64 machine, the C partitions took 350.60--372.36 seconds and the
Rust partitions at most 365.57 seconds. These timings are descriptive and are
not used in the proof.

The C and Rust programs deliberately implement the same orderly-generation
argument. Their exact agreement is strong evidence against language-specific,
state-management, and memory-safety errors, but it is not an independent
mathematical method or a third-party verification. The natural route to a
machine-checkable strengthening is a proof of nonexistence: encode
\emph{normalized detecting-matrix existence} as a Boolean satisfiability instance
and, on an \texttt{UNSAT} verdict, validate the emitted \textsc{drat} refutation
with an independent checker---or, after conversion to \textsc{lrat}, with a
formally verified checker such as \texttt{cake\_lpr} \cite{tan2021cakelpr}, as
done for the Boolean Pythagorean triples problem \cite{heule2016ptn}. This
succeeds on the small instances: we obtained checked refutations for every
nonexistence case through \(4\times6\) (hence \(\beta(Q_6)\geq5\)). It does not
reach the decisive scale. The instances are resolution-hard---a modern CDCL
solver did not close \(5\times7\) within minutes, whereas the enumerator settles
it in under a second---and the direct \(8\times14\) encoding already exceeds
\(6\times10^{7}\) clauses. Closing that size would require a substantially more
compact encoding together with cube-and-conquer partitioning, which we record as
the principal open direction for independent verification.

\begin{proof}[Proof of Theorem~\ref{thm:main}(i)]
The exhaustive computation finds no \(8\times14\) detecting matrix, so
Lemma~\ref{lem:equivalence} gives \(\beta(Q_{14})\geq9\). The explicit
\(9\times15\) detecting matrix in Appendix~\ref{app:certificates} yields, by
deleting any one of its columns, a \(9\times14\) detecting matrix. Hence
\(\beta(Q_{14})\leq9\).

If an \(8\times15\) detecting matrix existed, deleting any column would leave
an \(8\times14\) detecting matrix: a kernel vector for the smaller matrix could
be extended by a zero coordinate. Thus \(\beta(Q_{15})\geq9\), and the same
explicit \(9\times15\) matrix supplies the matching upper bound.
\end{proof}

\section{New upper bounds}\label{sec:upper-bounds}

The second part of Theorem~\ref{thm:main} is certified by the three sign
matrices reproduced in Appendix~\ref{app:certificates}. Their sizes and the
resulting bounds are summarized in Table~\ref{tab:bounds}.

\begin{table}[H]
\centering
\caption{New upper-bound certificates.}
\label{tab:bounds}
\begin{tabular}{cccc}
\toprule
Dimension & Rows & Previous upper bound & New upper bound\\
\midrule
24 & 13 & 14 & 13\\
26 & 14 & 15 & 14\\
29 & 15 & 16 & 15\\
\bottomrule
\end{tabular}
\end{table}

The previous values in the third column follow from the computational and
constructive bounds reported by Nikoli\'c et al., Hertz, and Lu--Ye
\cite{nikolic2017,hertz2020,luye2022}. Thus each certificate improves the
corresponding published bound by one.

\begin{proof}[Proof of Theorem~\ref{thm:main}(ii)]
For each matrix \(W\) in Appendix~\ref{app:certificates}, the exhaustive
meet-in-the-middle check of Section~\ref{sec:search} finds no nonzero
\(z\in\{-1,0,1\}^n\) satisfying \(Wz=0\). Hence each matrix is detecting, and
Lemma~\ref{lem:equivalence} gives the stated upper bound. Separate C and
JavaScript verifiers produce the same result -- one searching a sorted signature
table, the other a hash table, so agreement is not an artifact of shared code --
and both reject matrices with a deliberately duplicated column.
\end{proof}

The proof depends only on the displayed matrices and their finite
verification, not on the heuristic search that produced them. Deleting columns
also gives the corresponding bounds for all smaller dimensions, although those
consequences do not improve the previously known values at dimensions 23, 25,
27, and 28.

\subsection{Search methods}\label{sec:search}

The certificates were discovered in a separate upper-bound search built around a
single exact primitive: a meet-in-the-middle ternary-kernel counter. Splitting
the \(n\) columns into two halves, the routine enumerates the
\(3^{\lceil n/2\rceil}\) signed subset sums of one half, sorts their packed
row-sum signatures, and streams the signed subset sums of the other half,
counting exact cancellations \(Wz=0\). Each half contributes per-row sums of
small magnitude, so a full signature packs into two machine words with one
five-bit field per row and comparisons reduce to single instructions. This
replaces the \(3^n\) cost of naive ternary enumeration by
\(O\!\left(3^{\lceil n/2\rceil}\right)\) time and space; even for the
\(15\times29\) matrix (\(3^{15} = 14{,}348{,}907\) entries, occupying
\(\approx 115\)\,MB of memory), the kernel count is computed exactly in a
fraction of a second. The same primitive drives both the search and the final
check, so a discovered matrix is re-examined by an exact count rather than a
sampled estimate.

The \(15\times29\) matrix was found by simulated annealing on the \(\pm1\)
entries. To keep each proposal cheap, the annealing objective is a
\emph{bounded-weight} kernel count: only ternary vectors with at most a fixed
number of nonzero entries are enumerated. This underestimates the true count but
is far faster, and whenever it reaches zero the full exact counter re-certifies
the matrix; the exact null vectors it returns then steer the next move. Two
further reductions exploit the meet-in-the-middle split. Because the two halves
are independent, flipping one entry invalidates only one side's sorted table, so
each accepted move rebuilds a single \(3^{\lceil n/2\rceil}\)-sized side and
reuses the other unchanged. When the search stalls, repairs are restricted
to the coordinates in the support of the current null vectors, exploiting the
structure of the near-certificate rather than its count alone.

The remaining certificates were obtained without a fresh anneal. The
\(14\times26\) matrix is a row-and-column submatrix of the annealed
\(15\times29\) solution whose kernel count, evaluated once, is already zero. The
\(13\times24\) matrix came from an \emph{extension scan}: given a solved
\(13\times23\) submatrix, the two sorted half-tables are built once, and each of
the \(2^{13}\) candidate columns is then tested by a small update to one table
followed by a single merge, keeping exactly the columns whose addition creates no
ternary kernel vector. This reduces the cost of testing all \(2^{13}\) extending
columns from \(2^{13}\) independent kernel counts to a single table build
followed by \(2^{13}\) inexpensive merges; two independent \(13\times23\) slices
extended by the same column, corroborating the result. These procedures explain how the witnesses were found, but they are not
part of their proof.

\section{Coin weighing}

Let \(M(n)\) be the minimum number of fixed \(\bits\)-valued spring-scale
weighings that distinguish all subsets of \(n\) coins. The standard relation
\cite{sebotannier2004,luye2022} is
\[
 \beta(Q_n)-1\leq M(n)\leq\beta(Q_n).
\]
The left inequality follows by adjoining an all-ones row to a weighing matrix;
the right follows from normalizing a detecting matrix. It is only an
inequality, not an identity.

\begin{corollary}
\(M(14)=8\).
\end{corollary}

\begin{proof}
Theorem~\ref{thm:main} with the left inequality above gives \(M(14)\geq8\).
The artifact includes an
\(8\times14\) binary weighing matrix whose \(2^{14}\) subset sums are pairwise
distinct, giving \(M(14)\leq8\).
\end{proof}

\section{Reproducibility and limitations}

The repository separates fast checks from the decisive computation. The
command \texttt{./test.sh} verifies every upper-bound certificate with two
implementations, rejects broken controls, runs direct small-instance
comparisons, and exercises the Rust controls when the compiler is available.
The scripts \texttt{enumerate/run\_8x14\_audit.sh} and
\texttt{enumerate/run\_8x14\_rust\_audit.sh} execute all ten partitions and
write immutable-style bundles containing source, binary, environment, complete
logs, structural checks, and SHA-256 manifests.

For the next unresolved case, \(\beta(Q_{16})\), the artifact includes a portable
nine-row extension of the C enumerator and a deterministic frontier-sampling
driver, checked against all smaller controls and against the canonical form of
the known \(9\times15\) certificate. We do not report a runtime prediction,
because the work is highly uneven across frontier branches. In a bounded pilot
over the \(66{,}547\) canonical depth-six prefixes, the completed sample subtrees
ranged from 2 to \(30{,}964\) nodes while others remained unfinished at the time
limit; the driver therefore withholds a total-work estimate rather than
extrapolating from the easier units. Establishing nonexistence at \(n=16\) would
still require completing the entire frontier.

The public repository records the decisive node vectors, hashes, commands, and
validation results. \TODO{Before journal submission, the complete C and Rust
audit bundles should additionally be deposited in a durable public archive and
assigned a DOI. This is an artifact-availability task, not a missing step in
the enumeration itself, but it matters for long-term reproducibility.}

The literature search located prior uses of the underlying hypercube
symmetries, greedy and metaheuristic upper-bound searches, integer-programming
checks, and the coin-weighing bridge
\cite{mladenovic2012,nikolic2017,hertz2020,luye2022}. We found no prior report
of the exact values at dimensions 14 or 15, nor of this exhaustive
canonical-prefix computation. \TODO{Because absence claims are inherently
fragile, the novelty statement should be checked once more by a domain expert
and by the journal referee process.}

\section{Conclusion}

Normalizing one row converts the hypercube metric-dimension question into a
finite subset-sum enumeration with an explicit hypercube symmetry group.
Incremental collision testing makes each extension local, and orderly
generation under normalized row automorphisms removes noncanonical prefixes
without sacrificing completeness. The resulting exhaustive computation
establishes \(\beta(Q_{14})=\beta(Q_{15})=9\), while explicit certificates
improve the upper bounds for \(Q_{24}\), \(Q_{26}\), and \(Q_{29}\) and give
\(M(14)=8\). We also explored counterexample-guided abstraction
refinement (CEGAR), but the present formulation did not improve upon the
exhaustive method. Stronger abstractions, partition refinement (individualization-refinement),
and full row-pivoting symmetry reductions remain productive directions for scaling.
The natural next strengthening is an independently designed proof-producing
enumeration; the natural next open parameter is \(\beta(Q_{16})\).

\section*{Acknowledgments}
\addcontentsline{toc}{section}{Acknowledgments}

The author thanks Victor S. Miller for helpful discussions and for his
independent verification of the smaller terms of the sequence.

\section*{Disclosure of AI assistance}
\addcontentsline{toc}{section}{Disclosure of AI assistance}

The exploratory search code and the orchestration of the computational campaign
were developed with substantial assistance from large language models,
principally Anthropic's Claude and OpenAI's Codex; the same assistance was used
in drafting and revising this manuscript. The mathematical results do not rest
on that assistance. Every certificate and every mathematical nonexistence claim
is established by the standalone verifiers and the exhaustive enumerator, each
of which can be read, recompiled, and re-run independently of any AI system; the
heuristics that discovered the upper-bound matrices are separated from the
finite checks that prove them. The author has reviewed the definitions, proofs,
and claims and takes full responsibility for their correctness.

\clearpage
\appendix
\section{Detecting-matrix certificates}\label{app:certificates}

Each block below displays a sign matrix row by row, with \texttt{+} denoting
\(+1\) and \texttt{-} denoting \(-1\). By Lemma~\ref{lem:equivalence}, the
absence of a nonzero ternary kernel vector certifies the stated metric-dimension
upper bound. We include the \(9\times15\) matrix used in
Theorem~\ref{thm:main}(i) as well as the three new certificates from
Theorem~\ref{thm:main}(ii).

\noindent\begin{minipage}[t]{0.48\textwidth}
\vspace{0pt}
\subsection{A \texorpdfstring{\(9\times15\)}{9x15} certificate}
\begingroup\footnotesize
\begin{verbatim}
----++++++-++++
-++-+-++++-----
+-+++-+-----+++
++----++--+-+-+
-+-+---++--++--
+--------+---+-
++++-+-+-+-++-+
++++--+-++++++-
+-+-++--+--++--
\end{verbatim}
\endgroup

\subsection{A \texorpdfstring{\(13\times24\)}{13x24} certificate}
\begingroup\footnotesize
\begin{verbatim}
-+--+-++-+-++-+-+-+-++--
+-+--++++---+++++++--++-
--+-+---++-+-+-++-++-++-
+++--+-+++----++++++++++
-+-+-+--++-++-+--+-+-++-
+-++++-+++--++++--+++-+-
--++-+++++-+-+-+-+--+++-
++++++++++----++--+--++-
-+--+-++-++--+-+-+-+-++-
+-+--++++-++--+---+++-+-
--+-+---+++-+-++-+--+++-
+++--+-+++++++++--+--+++
+-++++-+++++--+++++--++-
\end{verbatim}
\endgroup

\end{minipage}\hfill
\begin{minipage}[t]{0.48\textwidth}
\vspace{0pt}
\subsection{A \texorpdfstring{\(14\times26\)}{14x26} certificate}
\begingroup\footnotesize
\begin{verbatim}
++++++++-+++++++++++++++++
--+--+-+-+-+-+-++-++-++-++
-+-+--+-++-+-++-+-+-++++++
++++---+++--+--++++++++++-
--+-+-+-++-+-+-+-+--+-++++
+--++--+++--++++--++++-+++
-+-+++-+++-+--+--+-+-+++++
++++++++++--+--+--+-++++++
--+--+-+-++-+-+--+--+-++++
+--+-++++-+++--+--++++-+++
-+-+--+-+++-++-+-+-+-+++++
++++---+++++++++--+-++++++
--+-+-+-+++-+-+-+-++-+++++
+--++--++++++--++++-++++++
\end{verbatim}
\endgroup

\subsection{A \texorpdfstring{\(15\times29\)}{15x29} certificate}
\begingroup\footnotesize
\begin{verbatim}
+++++++++-+++++++++++++++++++
--+--+-+-+-+-+-+-+-++-+-+-++-
+--+--+++++---+++++-+++--++++
-+-+-+--+-++-+--+-+++-++-++++
++++--+--+++--+--++++++++++++
--+-+-+-+-++-+-+-+-+-+-+-+-++
+--++++--+++--+++++---+++++-+
-+-++-++-+++-+--+-++-+--+-+++
++++++++++++--+--+++--+--++++
--+--+-+-+-++-+-+-++-+-+-+-++
+--+--+++++-+++--+-+--+++++-+
-+-+-+--+-+++-++-+++-+--+-+++
++++--+--+++++++++++--+--++++
--+-+-+-+-+++-+-+-+++-+-+-+++
+--++++--++++++--++++++--++++
\end{verbatim}
\endgroup
\end{minipage}

\bibliographystyle{plain}
\bibliography{references}
\end{document}
